\section*{Chapitre \ref{ch:pl:backtest-engine-architecture} --- Architecture du moteur de backtest}

\begin{solution}{exo:pl:backtest-engine-architecture:1}
L'annulation à $t = 2$ passe en premier~; puis, à $t = 3$, la mise à jour de la cotation (marché), l'exécution et l'ordre.
\end{solution}

\begin{solution}{exo:pl:backtest-engine-architecture:2}
Un backtest vectorisé n'a ni ordres, ni files, ni exécutions~: une stratégie de cotation lancée ici ne négocierait donc rien et renverrait zéro, un résultat qui
a l'air d'une réponse. Refuser rend visible le désaccord entre la stratégie et le niveau.
\end{solution}

\begin{solution}{exo:pl:backtest-engine-architecture:3}
Le niveau, un paramètre (latence, modèle de file, largeur de barre, limite d'inventaire) ou le hachage du code du moteur~: l'identifiant
hache l'ensemble du manifeste.
\end{solution}

\begin{solution}{exo:pl:backtest-engine-architecture:4}
Non~: $-12 \pm 13$ est à moins d'une erreur type de zéro. L'écart est le suivant~: le niveau~4 est inférieur au niveau~3 dans les vingt
sessions, de 90~\$ en moyenne avec une erreur type de 7.4, soit douze erreurs types de zéro.
\end{solution}

\begin{solution}{exo:pl:backtest-engine-architecture:5}
Une exécution partagée est la même transaction dans les deux passages~: même sens, même prix, même instant, contre le même ordre agressif de la session~; ce que fait
le prix dans les dix secondes suivantes relève surtout de la session, les markouts doivent donc concorder. Si ce n'était pas le cas, la classification
associerait des transactions différentes, ou les marchés des deux passages se seraient écartés avant les exécutions~; dans les deux cas, la
décomposition ne pourrait pas être lue.
\end{solution}

\begin{solution}{exo:pl:backtest-engine-architecture:6}
La réaction stratégique des autres participants~: les ordres à cours limité enregistrés ne sont pas recotés autour des ordres de la stratégie, les teneurs de marché
n'élargissent ni ne resserrent leur fourchette, et les traders informés ne modifient pas leur calendrier. Une population d'agents (\texttt{firm.agentmkt} du livre~10)
exécutée dans le simulateur l'ajouterait.
\end{solution}

\begin{solution}{exo:pl:backtest-engine-architecture:7}
Pour la graine~3 avec une limite de trois lots~: 417~\$, $-\$95.5$ et $-\$246$ aux niveaux~2, 3 et~4 (1\,462, 583 et 690 exécutions)~; avec dix lots~:
350~\$, $-\$110$ et $-\$472$ (1\,580, 670 et 818 exécutions). Plus d'inventaire coûte un peu aux niveaux~2 et~3 et beaucoup au niveau~4. La
part des exécutions réactives seules dans l'écart de markout passe de 129~\% à 147~\% (au-delà de 100~\% parce que, dans cette session, les exécutions du seul rejeu
perdent aussi)~: avec une limite plus grande, la stratégie continue de coter après avoir été dépassée, et au niveau~4 ces cotations sont réelles.
\end{solution}

\begin{solution}{exo:pl:backtest-engine-architecture:8}
Avec deux adaptations manuelles, l'écart entre les passages a trois causes possibles --- le modèle de marché, les adaptations et les
conventions des backtesters --- et l'affirmation en retient une sans preuve. Faire passer un même objet stratégie par un même moteur aux deux niveaux pour que les adaptations ne puissent pas
différer, puis comparer les journaux d'événements~: la première exécution divergente indique si le simulateur a traité différemment les deux cas, ou si la stratégie se trouvait simplement dans le carnet
(une exécution réactive seule). Un simulateur moins favorable parce que les ordres de la stratégie sont réels n'est pas un bogue.
\end{solution}

\begin{solution}{pb:pl:backtest-engine-architecture:1}
\begin{enumerate}
\item Une seule interface de stratégie, un seul \omterm{def:pl:backtest-engine-architecture:event}{modèle d'événements} et une seule horloge, une seule \omterm{def:pl:backtest-engine-architecture:dal}{couche d'accès aux données} et une seule \omterm{def:pl:backtest-engine-architecture:results}{base de résultats} pour l'ensemble des niveaux.
\item Une règle de cible pour le niveau~1 (sans ordres), exécutée barre par barre au niveau~2~; une stratégie de cotation avec des fonctions de rappel sur les marchés et les exécutions pour les niveaux~2
à~4.
\item L'ordre des événements simultanés~: données de marché, puis exécutions, puis ordres et annulations de la stratégie, de sorte que des entrées égales donnent des passages égaux et que
les niveaux soient comparables.
\item Les hachages du code du niveau, de la stratégie et du moteur, les hachages des données, les paramètres, la graine et les résultats synthétiques~; hacher l'ensemble fait que des entrées égales donnent un seul
identifiant, et tout changement un identifiant nouveau.
\item Pour que le passage enregistre les données qu'il a lues, quel que soit le nom des fichiers, et que «~mêmes données~» puisse être vérifié.
\item Le meilleur cours acheteur et le meilleur cours vendeur à la fin de chaque barre d'une seconde comportant des transactions~; une cotation s'exécute en totalité lorsque la barre se traite à son prix.
\item Une ombre prend une place dans la file enregistrée et s'exécute lorsque les exécutions enregistrées l'atteignent, sans modifier l'enregistrement~; un ordre de niveau~4
se pose dans la file du simulateur et se négocie contre les ordres rejoués, en modifiant ce qu'ils rencontrent ensuite.
\item Chaque ajout enregistré comme un ordre à cours limité, chaque annulation comme une annulation, chaque transaction comme un ordre au marché de sa taille.
\item La réaction stratégique des autres participants.
\item Par le simulateur qui tient lieu de plateforme, comme le livre~7 faisait intervenir un agent du marché synthétique en guise de substitut au trading réel.
\item Niveau~2~: 228~\$ ($\pm 29$) sur 841 exécutions~; niveau~3~: $-\$12$ ($\pm 13$) sur 333~; niveau~4~: $-\$102$ ($\pm 18$) sur 434.
\item Niveau~2 contre niveau~3, et niveau~3 contre niveau~4 (appariés par session)~; le niveau~3 ne se distingue pas de zéro.
\item Les vingt~; de 90~\$ par heure ($\pm 7.4$).
\item Dans leurs cinq premières exécutions de chaque session, à une médiane de 11~secondes dans l'heure.
\item Une exécution au contact n'exige qu'une transaction au prix, jamais une place dans une file.
\item \textbf{Résultat nommé.} P\&L par heure de 228~\$ au niveau~2, $-\$12$ au niveau~3 et $-\$102$ au niveau~4~; les exécutions réactives seules expliquent 101~\%
de l'écart de 68.4~\$ de markout à dix secondes entre le rejeu et la \omterm{def:pl:backtest-engine-architecture:reactive}{simulation réactive}, les exécutions partagées n'en expliquant aucune, et les
22~\$ restants sur les 90~\$ d'écart correspondent à un inventaire réévalué plus tard.
\item Les 241 exécutions réactives seules par session, qui perdent 69.2~\$ sur dix secondes~: des transactions qui n'ont lieu que parce que les ordres de la stratégie sont
réels, dans la file ou laissés au contact.
\item Environ 22~\$ par heure~: le P\&L de l'inventaire au-delà de dix secondes.
\item Qu'elle n'a pas d'edge~: elle équilibre au mieux derrière la file et perd là où ses ordres comptent~; le profit du backtest sur barres vient de l'exécution au
contact.
\item Elle fait de l'écart entre deux niveaux une propriété des marchés, puisque l'objet stratégie, le \omterm{def:pl:backtest-engine-architecture:event}{modèle d'événements} et les données sont les mêmes.
\end{enumerate}
\end{solution}

\begin{solution}{iq:pl:backtest-engine-architecture:1}
Aux exécutions~: combien, où et quand, à chaque niveau. Les backtests sur barres avec exécutions au contact ignorent la position dans la file et la sélection adverse~; la
simulation du carnet d'ordres ajoute les deux. Vérifier qu'il s'agit du même code de stratégie, puis décomposer l'écart par exécutions et par markouts.
\iqlookfor{Modèles d'exécution, files et sélection adverse~; une décomposition, non une supposition.}
\end{solution}

\begin{solution}{iq:pl:backtest-engine-architecture:2}
Par une règle fixe et documentée --- par exemple les données de marché avant les exécutions, puis les actions propres de la stratégie --- et ensuite par ordre d'arrivée, de sorte que
les passages soient déterministes et les niveaux comparables. Jamais par l'ordre d'une table de hachage.
\iqlookfor{Déterminisme par un ordre total.}
\end{solution}

\begin{solution}{iq:pl:backtest-engine-architecture:3}
Un manifeste par passage~: hachages du code de la stratégie et du moteur, hachages des instantanés de données, paramètres, graine et environnement~; un identifiant dérivé
de celui-ci~; des sorties rangées sous cet identifiant avec une interface de requête~; des données lues uniquement par une couche qui renvoie des hachages.
\iqlookfor{Identité du contenu du code et des données, non noms de fichiers ni dates.}
\end{solution}

\begin{solution}{iq:pl:backtest-engine-architecture:4}
Son impact propre~: ses ordres prennent des exécutions destinées à d'autres, restent visibles au contact et modifient ce que rencontrent les ordres suivants. Le faire tourner dans un
simulateur réactif sur la même session, puis décomposer l'écart en exécutions partagées, exécutions du seul rejeu et exécutions réactives seules, par
markout.
\iqlookfor{Ombres contre ordres réels~; un contrefactuel mesuré sur les mêmes données.}
\end{solution}

\begin{solution}{iq:pl:backtest-engine-architecture:5}
Deux formes de stratégie (une règle de cible et une stratégie pilotée par les événements) sur un même contexte~; des adaptateurs par backtester, les backtesters inchangés~; un
\omterm{def:pl:backtest-engine-architecture:event}{modèle d'événements} avec un ordre total~; une \omterm{def:pl:backtest-engine-architecture:dal}{couche d'accès aux données} à hachages de contenu~; une \omterm{def:pl:backtest-engine-architecture:results}{base de résultats} indexée par les hachages des manifestes~; des tests vérifiant que chaque adaptateur
reproduit le backtester appelé directement.
\iqlookfor{Adaptateurs sur les moteurs existants, un contexte, des passages reproductibles.}
\end{solution}
